
\subsubsection{2.1 Iterative Feedback Repair for LLM Code Generation}

SELF-REFINE (Madaan et al., 2023) introduces a framework where an LLM generates output, receives feedback, and iteratively refines. Applied to code, feedback takes the form of execution results, compiler messages, or test outputs. Olausson et al. (2023) provide an empirical survey of repair strategies for LLM-generated code, evaluating self-repair, external feedback, and execution-guided approaches across multiple benchmarks. Their findings indicate that repair effectiveness depends on the quality and specificity of the feedback signal.

These methods assume that the feedback signal (type errors, syntax errors, test failures) will be informative for the model-benchmark pair under study. This assumption is implicit and not empirically verified for specific model-benchmark combinations. The present work shows that for CodeLlama-7B on HumanEval, one class of formal feedback (mypy type checking) produces zero informative output across 2,680 samples.

Chen et al. (2022) introduce CodeT, which uses dual execution to filter LLM-generated solutions via test case agreement. CodeT demonstrates ensemble effectiveness but does not analyze failure mode complementarity between different repair channels.

\subsubsection{2.2 Grammar-Constrained Decoding}

Ugare et al. (2024) introduce SynCode, a grammar-constrained decoding framework that uses a pushdown automaton to mask syntactically invalid tokens during LLM generation. The present work integrates SynCode v0.4.16 into a multi-method pipeline and confirms a directional AST failure reduction (delta_ast = 0.075) on HumanEval with CodeLlama-7B. SynCode was previously evaluated in isolation without characterizing its interaction with post-hoc repair methods or measuring how it reshapes the downstream failure distribution.

\subsubsection{2.3 SMT-Based Verification and Repair}

Z3 (De Moura and Bjorner, 2008) is an SMT solver applicable to software verification. SMT-based repair approaches encode correctness constraints as logical formulae. Verifying or repairing code with Z3 requires problem-specific encodings: arithmetic constraints (e.g., integer-equality test assertions) are tractable, while general computation is not. No prior work has characterized Z3 applicability to standard code generation benchmarks. The present work quantifies this: 54 of 164 HumanEval problems (33\%) have integer-equality test assertions suitable for Z3 constraint encoding.

\subsubsection{2.4 Code LLM Evaluation}

HumanEval (Chen et al., 2021) is a standard benchmark for function-level Python code generation comprising 164 problems. Code Llama (Roziere et al., 2023) is the model family evaluated here. Roziere et al. report that smaller models exhibit higher syntax error rates, which is consistent with the finding that CodeLlama-7B failures are 97.5\% syntax-dominated. Existing evaluations report aggregate pass@k without decomposing failures by type. The FMD pipeline provides this decomposition.
